\begin{trapbox}
\begin{itemize}
    \item \textbf{左侧下界缺失}：很多同学只记得 $ab+bc+ca \leqslant a^2+b^2+c^2$，却忽略了当变量可以取负数时，$ab+bc+ca$ 存在下界 $-\frac{1}{2}(a^2+b^2+c^2)$。
    \item \textbf{等号条件}：必须满足 $a=b=c$。在约束条件下（如 $a,b,c$ 为三角形三边），等号可能取不到。
    \item \textbf{系数混淆}：注意分母是 $3$ 而非 $2$，这是三元对称结构最易记错的地方。
\end{itemize}
\end{trapbox}